\section{Requirements and Success Criteria}
\label{ch:requirements}

\begin{table}[htbp]
  \centering
  \caption{Traceable functional, non-functional and evidential requirements.}
  \label{tab:requirements}
  \small
  \begin{tabularx}{\textwidth}{p{0.11\textwidth}p{0.30\textwidth}Y}
    \toprule
    ID & Requirement & Evidence of success \\
    \midrule
    FR1 & Run autonomous agents in TORCS. & The runner exchanges telemetry and control actions with TORCS and records completed and failed episodes. \\
    FR2 & Support contrasting agent types. & Random, rule-based, map-aware, Dyna-Q and TD3 agents implement the common control contract. \\
    FR3 & Preserve experimental evidence. & SQLite records runs and telemetry; evaluation files retain policy path, trial outcome and contract information. \\
    FR4 & Provide a novice-facing learning workflow. & The GUI supports agent selection, live telemetry, results review, comparison and explanatory panels. \\
    NFR1 & Reduce deployment complexity. & The Windows package starts from \texttt{EnhancedAIRacing.exe} with TORCS and required runtime data included. \\
    NFR2 & Report reliability honestly. & Trial counts and failures remain visible; completed-lap medians are not presented as all-run averages. \\
    ER1 & Make behaviour interpretable. & Telemetry, outcome labels and agent-specific explanations link visible driving behaviour to AI concepts. \\
    ER2 & Support evidence-based learning. & The interface distinguishes capability, repeatability, failure and agent information sources rather than presenting a single high score. \\
    \bottomrule
  \end{tabularx}
\end{table}

The requirements are intentionally traceable to the research questions. FR1 and FR2 enable comparison; FR3 and NFR2 make the comparison defensible; FR4, NFR1, ER1 and ER2 address the brief's novice-access and learning objective. This traceability prevents the interface from becoming a feature list detached from the stated research contribution.

For the packaged release, the acceptance condition is practical: a user should not have to install Python, activate a virtual environment, configure a TORCS path or locate a policy checkpoint before observing a race. The release folder must remain intact because the executable depends on the adjacent internal runtime. The smoke test checks the novice route from opening the dashboard through starting TORCS to locating a completed run. This is more defensible than promising universal ease of use without a testable definition.

The educational requirements also constrain results presentation. The interface must not label an agent \enquote{best} solely because it has one fast record. It must show whether the claim is based on application history or repeated evaluation, include a trial count, and retain failures in the comparison. These requirements matter because an accessible wrapper that hides uncertainty would be easier to use but educationally misleading.

The requirements also define what the application deliberately does not do. It does not train a policy from the main beginner workflow, silently alter protected runtime checkpoints, or merge evidence from different evaluations into an unexplained aggregate. These exclusions preserve a stable interface and make the result of pressing Start understandable: the user is running a named, fixed agent under a selected scenario, not triggering an opaque background experiment.

